% GENERATED FILE. Edit canonical lessons, then run npm run generate:books.
% canonical-source-hash: fnv1a64:7af2bdb1b8e2f450
% canonical-lessons: GU-C332-dantaktha, GU-C332-bhagvan, GU-C332-daitya, GU-C332-devi, GU-C332-vardan

\chapter{Once Upon a Time}
\label{ch:gu-once-upon-a-time}
\hlchaptermodality{332}

\begin{chapteropening}
\textbf{By the end of this chapter:} \emph{I can retell the start of an old legend.}

Complete the last lesson of chapter 332: I can retell the start of an old legend.
\end{chapteropening}

\section[\texorpdfstring{\gu{દંતકથા} (dantakthā)}{દંતકથા (dantakthā)}]{\gu{દંતકથા} (dantakthā) — a legend}

\label{lesson:GU-C332-dantaktha}



\begin{quote}
Before the new one: say the Gujarati for a blessing, then the Gujarati for to offer (to a god).
\end{quote}

\subsection*{You'll want to know: \gu{દંતકથા}}
\textbf{\gu{દંતકથા}} — \emph{dantakthā} — \textquotedblleft{}a legend\textquotedblright{}.

\begin{quote}
\gu{દાદી} \gu{અમને} \gu{એક} \gu{જૂની} \gu{દંતકથા} \gu{કહેતાં} \gu{હતાં}.
\end{quote}

\begin{quote}
\emph{dādī amne ek jūnī dantakthā kahetā̃ hatā̃.}
\end{quote}

\begin{quote}
Grandmother used to tell us an old legend.
\end{quote}

A feminine noun.

\begin{cousinweb}
is built with \textbf{\gu{કથા}}, a tale, from Sanskrit \emph{kathā}, \textquotedblleft{}a story\textquotedblright{} (Monier-Williams); the history of its first part is not settled. A \textbf{\gu{દંતકથા}} is told from one generation to the next.
\end{cousinweb}

\begin{grammarlens}[title={What you've built}]
One more word for telling what happened.
\end{grammarlens}

\subsection*{Your turn}
\begin{itemize}
\raggedright
  \item \emph{Say it:} \emph{dantakthā}
  \item \emph{Say it:} \emph{dantakthā}, once more
  \item \emph{Say it:} \emph{dādī amne ek jūnī dantakthā kahetā̃ hatā̃.} — Grandmother used to tell us an old legend.
  \item \emph{Recall:} say the Gujarati for a blessing, then the Gujarati for to offer (to a god), then say \emph{dantakthā} again
\end{itemize}

\begin{tcolorbox}[breakable,colback=teal!4,colframe=teal!35!black,title={Before you move on}]
What does \gu{દંતકથા} mean? (\textquotedblleft{}A legend\textquotedblright{}.) Say it once more.
\end{tcolorbox}
\section[\texorpdfstring{\gu{ભગવાન} (bhagvān)}{ભગવાન (bhagvān)}]{\gu{ભગવાન} (bhagvān) — god; the Lord}

\label{lesson:GU-C332-bhagvan}



\begin{quote}
Before the new one: say the Gujarati for to offer (to a god), then the Gujarati for a legend.
\end{quote}

\subsection*{You'll want to know: \gu{ભગવાન}}
\textbf{\gu{ભગવાન}} — \emph{bhagvān} — \textquotedblleft{}god; the Lord\textquotedblright{}.

\begin{quote}
\gu{લોકોએ} \gu{ભગવાનને} \gu{પ્રાર્થના} \gu{કરી}.
\end{quote}

\begin{quote}
\emph{lokoe bhagvānne prārthanā karī.}
\end{quote}

\begin{quote}
The people prayed to god.
\end{quote}

A masculine noun.

\begin{cousinweb}
\textbf{\gu{ભગવાન}} (\emph{bhagvān}) is Sanskrit \emph{bhagavān}, the nominative of \emph{bhagavat}, \textquotedblleft{}fortunate, blessed\textquotedblright{}, and so \textquotedblleft{}the Lord\textquotedblright{} (Monier-Williams). Gujarati keeps the nominative ending as the whole word.
\end{cousinweb}

\begin{grammarlens}[title={What you've built}]
One more word for telling what happened.
\end{grammarlens}

\subsection*{Your turn}
\begin{itemize}
\raggedright
  \item \emph{Say it:} \emph{bhagvān}
  \item \emph{Say it:} \emph{bhagvān}, once more
  \item \emph{Say it:} \emph{lokoe bhagvānne prārthanā karī.} — The people prayed to god.
  \item \emph{Recall:} say the Gujarati for to offer (to a god), then the Gujarati for a legend, then say \emph{bhagvān} again
\end{itemize}

\begin{tcolorbox}[breakable,colback=teal!4,colframe=teal!35!black,title={Before you move on}]
What does \gu{ભગવાન} mean? (\textquotedblleft{}God; the Lord\textquotedblright{}.) Say it once more.
\end{tcolorbox}
\section[\texorpdfstring{\gu{દૈત્ય} (daitya)}{દૈત્ય (daitya)}]{\gu{દૈત્ય} (daitya) — a demon}

\label{lesson:GU-C332-daitya}



\begin{quote}
Before the new one: say the Gujarati for a legend, then the Gujarati for god.
\end{quote}

\subsection*{You'll want to know: \gu{દૈત્ય}}
\textbf{\gu{દૈત્ય}} — \emph{daitya} — \textquotedblleft{}a demon\textquotedblright{}.

\begin{quote}
\gu{જંગલમાં} \gu{એક} \gu{દૈત્ય} \gu{રહેતો} \gu{હતો}.
\end{quote}

\begin{quote}
\emph{jangalmā̃ ek daitya raheto hato.}
\end{quote}

\begin{quote}
A demon lived in the forest.
\end{quote}

A masculine noun.

\begin{cousinweb}
\textbf{\gu{દૈત્ય}} (\emph{daitya}) is a tatsama, Sanskrit \emph{daitya}, \textquotedblleft{}a son of Diti, a demon\textquotedblright{} (Monier-Williams): in the old stories, the enemies of the gods.
\end{cousinweb}

\begin{grammarlens}[title={What you've built}]
One more word for telling what happened.
\end{grammarlens}

\subsection*{Your turn}
\begin{itemize}
\raggedright
  \item \emph{Say it:} \emph{daitya}
  \item \emph{Say it:} \emph{daitya}, once more
  \item \emph{Say it:} \emph{jangalmā̃ ek daitya raheto hato.} — A demon lived in the forest.
  \item \emph{Recall:} say the Gujarati for a legend, then the Gujarati for god, then say \emph{daitya} again
\end{itemize}

\begin{tcolorbox}[breakable,colback=teal!4,colframe=teal!35!black,title={Before you move on}]
What does \gu{દૈત્ય} mean? (\textquotedblleft{}A demon\textquotedblright{}.) Say it once more.
\end{tcolorbox}
\section[\texorpdfstring{\gu{દેવી} (devī)}{દેવી (devī)}]{\gu{દેવી} (devī) — a goddess}

\label{lesson:GU-C332-devi}



\begin{quote}
Before the new one: say the Gujarati for god, then the Gujarati for a demon.
\end{quote}

\subsection*{You'll want to know: \gu{દેવી}}
\textbf{\gu{દેવી}} — \emph{devī} — \textquotedblleft{}a goddess\textquotedblright{}.

\begin{quote}
\gu{દેવીએ} \gu{દૈત્યને} \gu{હરાવ્યો}.
\end{quote}

\begin{quote}
\emph{devīe daityane harāvyo.}
\end{quote}

\begin{quote}
The goddess defeated the demon.
\end{quote}

A feminine noun.

\begin{cousinweb}
\textbf{\gu{દેવી}} (\emph{devī}) is a tatsama, Sanskrit \emph{devī}, \textquotedblleft{}a goddess\textquotedblright{} (Monier-Williams). As the doer she takes \textbf{\gu{એ}}, and the verb agrees with the demon defeated: \textbf{\gu{દૈત્યને} \gu{હરાવ્યો}}.
\end{cousinweb}

\begin{grammarlens}[title={What you've built}]
One more word for telling what happened.
\end{grammarlens}

\subsection*{Your turn}
\begin{itemize}
\raggedright
  \item \emph{Say it:} \emph{devī}
  \item \emph{Say it:} \emph{devī}, once more
  \item \emph{Say it:} \emph{devīe daityane harāvyo.} — The goddess defeated the demon.
  \item \emph{Recall:} say the Gujarati for god, then the Gujarati for a demon, then say \emph{devī} again
\end{itemize}

\begin{tcolorbox}[breakable,colback=teal!4,colframe=teal!35!black,title={Before you move on}]
What does \gu{દેવી} mean? (\textquotedblleft{}A goddess\textquotedblright{}.) Say it once more.
\end{tcolorbox}
\section[\texorpdfstring{\gu{વરદાન} (vardān)}{વરદાન (vardān)}]{\gu{વરદાન} (vardān) — a boon}

\label{lesson:GU-C332-vardan}



\begin{quote}
Before the new one: say the Gujarati for a demon, then the Gujarati for a goddess.
\end{quote}

\subsection*{You'll want to know: \gu{વરદાન}}
\textbf{\gu{વરદાન}} — \emph{vardān} — \textquotedblleft{}a boon\textquotedblright{}.

\begin{quote}
\gu{ભગવાને} \gu{રાજાને} \gu{વરદાન} \gu{આપ્યું}.
\end{quote}

\begin{quote}
\emph{bhagvāne rājāne vardān āpyũ.}
\end{quote}

\begin{quote}
The god granted the king a boon.
\end{quote}

A neuter noun.

\begin{cousinweb}
\textbf{\gu{વરદાન}} (\emph{vardān}) is Sanskrit \emph{varadāna}, \textquotedblleft{}the granting of a boon\textquotedblright{} (Monier-Williams), built from \emph{vara}, a wish chosen, and \emph{dāna}, a giving. A god \textbf{\gu{વરદાન} \gu{આપે} \gu{છે}}.
\end{cousinweb}

\begin{grammarlens}[title={What you've built}]
One more word for telling what happened.
\end{grammarlens}

\subsection*{Your turn}
\begin{itemize}
\raggedright
  \item \emph{Say it:} \emph{vardān}
  \item \emph{Say it:} \emph{vardān}, once more
  \item \emph{Say it:} \emph{bhagvāne rājāne vardān āpyũ.} — The god granted the king a boon.
  \item \emph{Recall:} say the Gujarati for a demon, then the Gujarati for a goddess, then say \emph{vardān} again
\end{itemize}

\begin{tcolorbox}[breakable,colback=teal!4,colframe=teal!35!black,title={Before you move on}]
What does \gu{વરદાન} mean? (\textquotedblleft{}A boon\textquotedblright{}.) Say it once more.
\end{tcolorbox}
